\chapter{Mathematical Foundations}
\label{chap:math}

\section{Coordinates and Constants}

The vertical coordinate is the \textbf{Exner pressure}
\begin{equation}
\boxed{\;\pi \equiv \left(\frac{p}{p_0}\right)^{\!\kappa}\;,\qquad
       p_0=10^{5}\,\mathrm{Pa},\qquad
       \kappa=\frac{R_d}{c_p}=\tfrac{2}{7}\approx 0.2857\;}
\label{eq:exner}
\end{equation}
with $\pi=1$ at the surface and $\pi\to 0$ at the model top.  The Wu
Fortran hardcodes $\kappa=2/7$ (this exact value is what makes the
$Q\to$\,SI conversion a clean constant, \Cref{chap:wu_coeffs}).  On the
$N_L=9$ ERA5 pressure levels used here,
\begin{equation}
\mathrm{PLEVS}=[1000,850,700,500,400,300,250,200,100]\ \mathrm{hPa},
\end{equation}
\begin{equation}
\pi_k = [\,1.000,\,0.955,\,0.903,\,0.820,\,0.770,\,0.709,\,0.673,\,0.631,\,0.518\,].
\end{equation}

\begin{center}\small
\begin{tabular}{lll}
  \toprule
  Constant & Symbol & Value \\
  \midrule
  Gravity (Fortran)      & $g$        & $9.81$ m\,s$^{-2}$ \\
  Dry-air gas constant   & $R_d$      & $287$ J\,kg$^{-1}$\,K$^{-1}$ \\
  Specific heat at $p$   & $c_p$      & $1004.5$ J\,kg$^{-1}$\,K$^{-1}$ (Fortran)\\
  $\kappa=R_d/c_p$       & $\kappa$   & $2/7 \approx 0.2857$ \\
  Earth's rotation       & $\Omega$   & $7.292\times10^{-5}$ rad\,s$^{-1}$ \\
  Earth radius (Wu)      & $a=2{\times}10^{7}/\pi$ & $6.366\times10^{6}$ m  \\
  Reference pressure     & $p_0$      & $10^{5}$ Pa \\
  Reference Coriolis     & $f_0$      & $f(45^\circ)\approx 1.031\times10^{-4}$ s$^{-1}$ \\
  \bottomrule
\end{tabular}
\end{center}

\section{Thermodynamic Relations}

\textbf{Potential temperature}
\begin{equation}
\theta = T\left(\frac{p_0}{p}\right)^{\!\kappa}.
\end{equation}
The Fortran is fed \emph{raw temperature} $T$ and forms $\theta$
internally as $\theta = T\,c_p/\Pi$ with $\Pi=c_p(p/p_0)^\kappa$
(\texttt{pvpialln\_94UV.f}); feeding it $\theta$ would double-convert and
inflate the static stability.

\textbf{Hydrostatic balance in Exner coordinates}
\begin{equation}
\boxed{\;\frac{\partial\Phi}{\partial\pi}=-c_p\,\theta\;}
\label{eq:hydrostatic}
\end{equation}
is the vertical link between $\Phi$ and $\theta$ that closes the
inversion and \emph{defines the boundary condition} of
\Cref{sec:boundary_theta}.

\section{Ertel Potential Vorticity}

The hydrostatic Ertel PV is
\begin{equation}
Q = -g(\zeta+f)\frac{\partial\theta}{\partial p}
    - g\Bigl(\frac{\partial u}{\partial p}\frac{\partial\theta}{\partial y}
            - \frac{\partial v}{\partial p}\frac{\partial\theta}{\partial x}\Bigr),
\label{eq:Q}
\end{equation}
in SI units K\,m$^2$\,kg$^{-1}$\,s$^{-1}$ ($\equiv 10^{-6}$\,PVU).  The
Wu Fortran evaluates an Exner-coordinate form of \cref{eq:Q}; its exact
relation to SI is derived in \Cref{chap:wu_coeffs}.

\begin{warnbox}[Interior PV exists only on the interior levels]
\texttt{pvpialln} computes $Q$ only on the \emph{interior} levels
$K=2\ldots N_L-1$ (850--200\,hPa): the vertical derivative
$\partial\theta/\partial\pi$ needs neighbours on both sides.  At the
surface ($K{=}1$, 1000\,hPa) and top ($K{=}N_L$, 100\,hPa) there is
\textbf{no interior PV} --- those two levels instead carry the
\emph{boundary potential temperature} of \Cref{sec:boundary_theta}.
This is why a surface-confined thermal anomaly shows ``no PV at
1000\,hPa'' in the interior field.
\end{warnbox}

\section{Boundary Potential Temperature as a PV Sheet}
\label{sec:boundary_theta}

The two boundary levels are represented by the half-level-averaged
potential temperature (\texttt{pvpialln\_94UV.f}, the lines that set
\texttt{THB,THT}):
\begin{equation}
\theta_b^{\text{(bot)}} = \tfrac12\bigl(\theta_{K=1}+\theta_{K=2}\bigr),\qquad
\theta_b^{\text{(top)}} = \tfrac12\bigl(\theta_{K=N_L-1}+\theta_{K=N_L}\bigr).
\end{equation}
In the \texttt{q.out} file the layout is therefore
$[\,\theta_b^{\text{bot}},\ \theta_b^{\text{top}},\ Q_{K=2},\ldots,Q_{K=N_L-1}\,]$
--- boundary $\theta$ first, interior PV after.

Through hydrostatic balance \cref{eq:hydrostatic}, $\theta_b$ specifies
the \emph{vertical derivative} of the geopotential / streamfunction at
the boundary, i.e.\ an \textbf{inhomogeneous Neumann boundary
condition}.  In \texttt{qinvert21\_94.f} (BALNC) it enters the elliptic
stencil one node inside each boundary (the lines that add
$+\theta_b^{\text{bot}}/\Delta\pi$ at $K{=}2$ and
$-\theta_b^{\text{top}}/\Delta\pi$ at $K{=}N_L-1$) and reconstructs the
boundary layer by a half-step hydrostatic integration,
$H_{1}=H_{2}+\theta_b^{\text{bot}}(\pi_2-\pi_1)$.  If no boundary
$\theta$ is supplied, the code applies \emph{homogeneous} Neumann
conditions.

\begin{keybox}[Bretherton: surface theta as a boundary PV sheet]
A warm surface potential-temperature anomaly is dynamically equivalent
to a \emph{positive} PV anomaly concentrated as a delta-function sheet
at the lower boundary; a cold anomaly $\equiv$ a negative sheet.  By the
PV--inversion (``invertibility'') principle this induces a circulation:
\[
\theta_b'>0\ (\text{warm}) \;\Rightarrow\; \text{cyclonic } \psi'\ (\psi'<0 \text{ in NH}),
\qquad
\theta_b'<0\ (\text{cold}) \;\Rightarrow\; \text{anticyclonic}.
\]
This is the textbook ``PV-thinking'' interpretation
\cite{HoskinsMcIntyreRobertson1985} of the Neumann condition above; the
Fortran implements the boundary condition, not the interpretation.  It
is the reason the \textbf{surface piece is inverted separately}
(its PV ``lives in'' the boundary $\theta$), and why for this event the
surface-piece $\psi'$ is predominantly cyclonic (the domain-mean surface
$\theta'$ is net warm, $\approx +1$\,K).
\end{keybox}

\section{Linearization and Balance Closure}
\label{sec:closure}

Decompose $\psi=\bpsi+\dpsi$, $\Phi=\bPhi+\dPhi$.  The piecewise solver
inverts the PV/boundary-$\theta$ \emph{perturbation}
$\dq=q_{\text{event}}-q_{\text{mean}}$ under a balance constraint.  The
production configuration uses the \textbf{nonlinear} balance
(\texttt{INLIN}=1, \Cref{sec:wu_params}): the operator coefficients are
evaluated about the half-state $\bar\cdot + \tfrac12(\cdot)'$
(\texttt{CNLIN}=0.5 in \texttt{qinvertp21\_94.f}), so that the three
piecewise solutions \emph{sum to the total} perturbation,
\begin{equation}
\boxed{\;\dpsi_{\text{surface}}+\dpsi_{\text{lower}}+\dpsi_{\text{upper}}
       \;\approx\;\dpsi_{\text{total}}\;}
\label{eq:additivity}
\end{equation}
to solver tolerance (the linear-closure check of \Cref{sec:step11}).

\begin{warnbox}[Why INLIN=1 rather than INLIN=0]
Earlier notes claimed \texttt{INLIN}=0 (linear balance) was
\emph{mandatory} because the nonlinear balance ``explodes at the jet.''
On the current 9-level, extended-domain configuration the
\textbf{nonlinear balance \texttt{INLIN}=1 converges} (all three pieces
reach \texttt{TOTAL CONVERGENCE}) and is required for the additivity
\cref{eq:additivity}.  \texttt{INLIN}=0 is no longer used.
\end{warnbox}
